% Part: many-valued-logic
% Chapter: three-valued-logics
% Section: kleene

\documentclass[../../../include/open-logic-section]{subfiles}

\begin{document}

\olfileid[pt]{mvl}{thr}{skl}

\olsection{Lógicas de Kleene}

Stephen Kleene introduziu duas lógicas trivaloradas motivadas por uma lógica
na qual os valores de verdade são concebidos como resultados de procedimentos
computacionais: um procedimento pode produzir $\True$ ou $\False$, mas também pode
não terminar. Nesse caso, o valor de verdade correspondente é
indefinido, representado pelo valor de verdade~$\Undef$.

Para calcular a negação de uma proposição~$!A$, primeiro se calcularia
o valor de~$!A$ e, em seguida, seria devolvido o oposto do resultado. Se o
cálculo de $!A$ não termina, o procedimento inteiro também não
termina: portanto, a negação de $\Undef$ é $\Undef$.

Para calcular uma conjunção $!A \land !B$, há duas opções: pode-se
calcular primeiro~$!A$, depois $!B$, e então o resultado será $\True$ se
o resultado de ambos for~$\True$, e $\False$ caso contrário. Se qualquer um dos
cálculos não parar, o procedimento inteiro tampouco para. Assim,
nesse caso, se um dos membros da conjunção for indefinido, a conjunção também será.
O mesmo vale para a disjunção.

Contudo, se pudermos avaliar $!A$ e $!B$ em paralelo, poderemos fazer melhor.
Então, se um dos dois procedimentos parar e devolver $\False$, poderemos
parar, pois a resposta deve ser falsa. Nesse caso, portanto, uma conjunção com
um membro falso é falsa, mesmo que o outro seja indefinido.
De modo semelhante, ao calcular uma disjunção em paralelo, podemos parar assim que
o procedimento de um dos dois disjuntos devolver verdadeiro: a
disjunção deve então ser verdadeira. Nesse caso, podemos saber qual é o resultado
de uma afirmação composta, mesmo que um de seus componentes seja indefinido. Sob
essa interpretação, poderíamos ler $\Undef$ como ``desconhecido'', e não como
``indefinido''.

As duas interpretações dão origem às lógicas forte e fraca de Kleene.
O condicional é definido como equivalente a $\lnot !A \lor !B$.

\begin{defn}
A \emph{lógica forte de Kleene}~$\LogKs$ é definida pela matriz:
\begin{enumerate}
  \item A linguagem proposicional padrão $\Lang L_0$, com
  $\lnot$, $\land$, $\lor$, $\lif$.
  \item O conjunto de valores de verdade $V = \{\True, \Undef, \False\}$.
  \item $\True$ é o único valor designado, isto é, $V^+ = \{\True\}$.
  \item As funções de verdade são dadas pelas seguintes tabelas:
  \begin{center}
    \begin{tabular}{c|c} 
      $\tf{\lnot}$ & \\ 
      \hline  
      $\True$ & $\False$ \\ 
      $\Undef$ & $\Undef$ \\
      $\False$ & $\True$ 
    \end{tabular}
    \quad
    \begin{tabular}{c|ccc} 
      $\tf{\land}[\LogKs]$ & $\True$ & $\Undef$ & $\False$ \\ 
      \hline 
      $\True$ & $\True$ & $\Undef$ & $\False$ \\ 
      $\Undef$ & $\Undef$ & $\Undef$ & $\False$\\ 
      $\False$ & $\False$ & $\False$ & $\False$ 
    \end{tabular}
    \\[2ex]
    \begin{tabular}{c|ccc} 
      $\tf{\lor}[\LogKs]$ & $\True$ & $\Undef$ & $\False$ \\ 
      \hline 
      $\True$ & $\True$ & $\True$ & $\True$ \\ 
      $\Undef$ & $\True$ & $\Undef$ & $\Undef$ \\
      $\False$ & $\True$ & $\Undef$ & $\False$ 
    \end{tabular}
    \quad
    \begin{tabular}{c|ccc} 
      $\tf{\lif}[\LogKs]$ & $\True$ & $\Undef$ & $\False$ \\ 
      \hline 
      $\True$ & $\True$ & $\Undef$ & $\False$ \\ 
      $\Undef$ & $\True$ & $\Undef$ & $\Undef$  \\ 
      $\False$ & $\True$ & $\True$ & $\True$ 
    \end{tabular}
  \end{center} 
\end{enumerate}
\end{defn}

\begin{defn}
A \emph{lógica fraca de Kleene}~$\LogKw$ é definida pela matriz:
\begin{enumerate}
  \item A linguagem proposicional padrão $\Lang L_0$, com
  $\lnot$, $\land$, $\lor$, $\lif$.
  \item O conjunto de valores de verdade $V = \{\True, \Undef, \False\}$.
  \item $\True$ é o único valor designado, isto é, $V^+ = \{\True\}$.
  \item As funções de verdade são dadas pelas seguintes tabelas:
  \begin{center}
    \begin{tabular}{c|c} 
      $\tf{\lnot}$ & \\ 
      \hline  
      $\True$ & $\False$ \\ 
      $\Undef$ & $\Undef$ \\
      $\False$ & $\True$ 
    \end{tabular}
    \quad
    \begin{tabular}{c|ccc} 
      $\tf{\land}[\LogKw]$ & $\True$ & $\Undef$ & $\False$ \\ 
      \hline 
      $\True$ & $\True$ & $\Undef$ & $\False$ \\ 
      $\Undef$ & $\Undef$ & $\Undef$ & $\Undef$\\ 
      $\False$ & $\False$ & $\Undef$ & $\False$ 
    \end{tabular}
    \\[2ex]
    \begin{tabular}{c|ccc} 
      $\tf{\lor}[\LogKw]$ & $\True$ & $\Undef$ & $\False$ \\ 
      \hline 
      $\True$ & $\True$ & $\Undef$ & $\True$ \\ 
      $\Undef$ & $\Undef$ & $\Undef$ & $\Undef$ \\
      $\False$ & $\True$ & $\Undef$ & $\False$ 
    \end{tabular}
    \quad
    \begin{tabular}{c|ccc} 
      $\tf{\lif}[\LogKw]$ & $\True$ & $\Undef$ & $\False$ \\ 
      \hline 
      $\True$ & $\True$ & $\Undef$ & $\False$ \\ 
      $\Undef$ & $\Undef$ & $\Undef$ & $\Undef$  \\ 
      $\False$ & $\True$ & $\Undef$ & $\True$ 
    \end{tabular}
  \end{center} 
\end{enumerate}
\end{defn}

\begin{prop}
  $\LogKs$ e $\LogKw$ não têm tautologias.
\end{prop}

\begin{proof}
  Se $\pAssign v(p) = \Undef$ para todas as !!{propositional variable}s~$p$,
  então qualquer fórmula~$!A$ terá o valor de verdade~$\pValue v(!A) =
  \Undef$, pois
  \[
    \tf{\lnot}(\Undef) = \tf{\lor}(\Undef, \Undef) = \tf{\land}(\Undef,
  \Undef) = \tf{\lif}(\Undef, \Undef) = \Undef
  \]
  em ambas as lógicas. Como $\Undef \notin V^+$ tanto para $\LogKs$ quanto para
  $\LogKw$, nessa !!{valuation}, $!A$ não será designada.
\end{proof}

Embora as lógicas fraca e forte de Kleene não tenham tautologias, elas
têm relações de consequência não triviais. 

\begin{prob}
  Quais das relações a seguir valem (a) na lógica forte e
  (b) na lógica fraca de Kleene? Dê uma tabela-verdade para cada uma.
  \begin{enumerate}
    \item $p, p \lif q \Entails q$
    \item $p \lor q, \lnot p \Entails q$
    \item $p \land q \Entails p$
    \item $p \Entails p \land p$
    \item $p \Entails p \lor q$
  \end{enumerate}
\end{prob}

Dmitry Bochvar interpretou $\Undef$ como ``sem sentido'' e tentou
usá-lo para resolver paradoxos como o paradoxo do Mentiroso, estipulando
que sentenças paradoxais assumem o valor~$\Undef$. Ele introduziu uma
lógica que é, em essência, a lógica fraca de Kleene ampliada por conectivos
adicionais, dois dos quais são a ``negação externa'' e o operador ``é
indefinido'':
\begin{center}
  \begin{tabular}{c|c} 
    $\tf{\sim}$ & \\ 
    \hline  
    $\True$ & $\False$ \\ 
    $\Undef$ & $\True$ \\
    $\False$ & $\True$ 
  \end{tabular}
\quad
  \begin{tabular}{c|c} 
    $\tf{+}$ & \\ 
    \hline  
    $\True$ & $\False$ \\ 
    $\Undef$ & $\True$ \\
    $\False$ & $\False$ 
  \end{tabular}
\end{center}

\begin{prob}
  Você consegue definir $\sim$ na lógica de Bochvar em termos de $\lnot$
  e~$+$, isto é, encontrar !!a{formula} que contenha apenas a !!{propositional
  variable}~$p$ e não envolva $\sim$, mas sempre assuma o mesmo
  valor de verdade que~$\mathord{\sim}p$? Dê uma tabela-verdade que demonstre
  sua resposta.
\end{prob}


\end{document}
